\documentclass[11pt,letterpaper]{article}
\usepackage[top=0.7in,bottom=0.7in,left=0.75in,right=0.75in]{geometry}
\usepackage[T1]{fontenc}
\usepackage[utf8]{inputenc}
\usepackage{lmodern}
\usepackage{microtype}
\usepackage{enumitem}
\usepackage[hidelinks]{hyperref}
\usepackage{titlesec}
\setlength{\parindent}{0pt}
\setlength{\parskip}{0pt}
\pagestyle{empty}
\titleformat{\section}{\normalfont\bfseries}{}{0em}{\uppercase}[\vspace{-4pt}\hrule\vspace{4pt}]
\titlespacing{\section}{0pt}{8pt}{5pt}
\setlist[itemize]{leftmargin=1.2em,label=--,itemsep=0pt,topsep=2pt,parsep=0pt}
\begin{document}
\begin{center}
{\fontsize{16}{19}\selectfont\textbf{\MakeUppercase{Diego Castillo Pineda}}}\\[4pt]
{\small La Florida, Santiago, Chile \quad|\quad +56 9 9416 2547 \quad|\quad \href{mailto:diego.castillop11@gmail.com}{diego.castillop11@gmail.com}\\
\href{https://linkedin.com/in/diego-castillo-pineda}{linkedin.com/in/diego-castillo-pineda} \quad|\quad \href{https://github.com/diegocastillop}{github.com/diegocastillop}}
\end{center}

\section{Resumen Profesional}
Ingeniero de Datos con sólida experiencia en administración y optimización de bases de datos SQL Server, Azure SQL y Oracle en entornos de alta concurrencia transaccional. Especializado en automatización de procesos mediante Python y T-SQL, diseño de pipelines de datos y desarrollo de soluciones analíticas con Power BI para apoyar decisiones de negocio. Experiencia comprobada en migración cloud hacia Microsoft Azure, integración de datos en formatos JSON/XML y optimización de consultas que redujeron tiempos de respuesta en hasta 60\% durante eventos críticos con más de 50,000 transacciones por hora.

\section{Experiencia Profesional}

\textbf{Punto Ticket} \hfill Santiago, Chile\\
\textit{Analista de Base de Datos Senior} \hfill Nov. 2019 – Feb. 2026
\begin{itemize}
  \item Diseñé y optimicé stored procedures en SQL Server para procesar más de 50,000 transacciones concurrentes por hora durante eventos masivos, manteniendo tiempos de respuesta bajo 2 segundos según SLA operativo.
  \item Automaticé procesos de carga masiva y transformación de datos mediante scripts T-SQL y Python, reduciendo errores manuales en 85\% y liberando 15 horas semanales del equipo para tareas analíticas.
  \item Desarrollé consultas dinámicas con PIVOT y agregaciones complejas para reportería ejecutiva, reemplazando procesos manuales en Excel y entregando dashboards automatizados en Power BI.
  \item Integré sistemas mediante parseo y transformación de datos en JSON, XML y HTML, asegurando la consistencia entre plataformas transaccionales y analíticas.
  \item Participé activamente en la migración de bases de datos on-premise hacia Azure SQL Database, colaborando con equipos DevOps para implementar estrategias de disaster recovery y alta disponibilidad.
  \item Administré instancias SQL Server en entornos productivos críticos, monitoreando performance con DMVs, ejecutando tuning de índices y planes de ejecución para garantizar disponibilidad 99.9\%.
\end{itemize}
\vspace{2pt}
\textbf{Imperial S.A.} \hfill Santiago, Chile\\
\textit{Analista Funcional} \hfill Jun. 2017 – Nov. 2019
\begin{itemize}
  \item Levanté requerimientos funcionales junto a usuarios clave de finanzas y operaciones, traduciéndolos en especificaciones técnicas para desarrollo en Oracle y SQL Server.
  \item Desarrollé consultas SQL complejas en Oracle y SQL Server Management Studio para extracción, transformación y análisis de datos de sistemas ERP, apoyando decisiones operativas y auditorías.
  \item Parametricé módulos en sistemas ERP (Oracle) según especificaciones funcionales, asegurando la correcta implementación de reglas de negocio y validaciones de datos.
  \item Elaboré documentación técnica y funcional (manuales de usuario, diagramas de flujo, instructivos) que mejoraron la transferencia de conocimiento y reducción de tickets de soporte en 30\%.
  \item Brindé soporte directo a usuarios finales en ambientes productivos, diagnosticando incidencias relacionadas con integridad de datos, performance de consultas y configuración de sistemas.
\end{itemize}
\vspace{2pt}
\textbf{OB GROUP PARK} \hfill Santiago, Chile\\
\textit{Analista de Sistemas y TI} \hfill Mayo 2014 – Abril 2017
\begin{itemize}
  \item Lideré proyectos de integración tecnológica entre sistemas contables y plataformas retail, coordinando con equipos multidisciplinarios para asegurar compatibilidad de datos y continuidad operativa.
  \item Desarrollé scripts de automatización para conciliación de datos entre sistemas transaccionales y contables, reduciendo tiempos de cierre mensual en 40\%.
  \item Capacité a usuarios finales en el uso de nuevas herramientas tecnológicas implementadas, elaborando material didáctico y sesiones prácticas que mejoraron la adopción en 90\%.
\end{itemize}
\vspace{2pt}
\textbf{Hewlett-Packard Company} \hfill Santiago, Chile\\
\textit{Agente de Mesa de Ayuda Nivel 1} \hfill Ago. 2013 – Mayo 2014
\begin{itemize}
  \item Brindé soporte técnico a usuarios finales, diagnosticando y resolviendo incidencias de software, hardware y acceso a sistemas, cumpliendo SLA de resolución en primer contacto del 75\%.
\end{itemize}

\section{Proyectos Destacados}

\textbf{Sistema Inteligente de Gestión de Cobranzas con IA} \hfill 2024
\begin{itemize}
  \item Desarrollé una aplicación full-stack con React, Node.js y PostgreSQL que utiliza prompt engineering con Claude AI para clasificar automáticamente gestiones de cobranza y priorizar casos críticos según patrones de comportamiento de pago.
  \item Implementé pipeline ETL en Python para procesar datos históricos de cobranzas, generar métricas predictivas de recuperación y visualizarlas en dashboards Power BI interactivos.
  \item Diseñé arquitectura de base de datos segura con encriptación de datos sensibles, control de acceso basado en roles y auditoría de transacciones, cumpliendo estándares de protección de datos financieros.
  \item Apliqué principios de UI/UX para crear interfaces intuitivas que redujeron el tiempo promedio de registro de gestiones en 50\%, validado con usuarios reales del sector financiero.
\end{itemize}

\section{Habilidades T\'{e}cnicas}
\textbf{Datos \& Cloud:} SQL Server (Experto), Azure SQL Database, Oracle, PostgreSQL, MongoDB, Microsoft Azure (Azure Data Factory, Azure Functions), ETL/ELT \\
\textbf{Desarrollo \& IA:} Python (pandas, scikit-learn), C\#, Node.js, Go, Prompt Engineering (Claude AI, Gemini), Machine Learning básico \\
\textbf{Frontend \& Diseño:} React, TypeScript, JavaScript (ES6+), HTML5/CSS3, UI/UX Design, Responsive Design \\
\textbf{Herramientas \& Metodologías:} Git (GitHub/GitLab), Docker, Power BI (DAX, Power Query), Excel Avanzado (VBA, Power Pivot), T-SQL, Visual Studio Code, Scrum/Kanban, Jira

\section{Formaci\'{o}n Acad\'{e}mica}
\textbf{Diplomado en Business Analytics \& Data Science}, Universidad de Chile \hfill 2023 \\
\textbf{Especialización en Machine Learning con Python}, Universidad de Chile \hfill 2023 \\
\textbf{Especialización en Administración de Bases de Datos SQL Server}, Universidad de Chile \hfill 2023 \\
\textbf{Desarrollo de Aplicaciones Móviles}, Universidad Complutense de Madrid \hfill 2017 \\
\textbf{Analista Programador}, Universidad Tecnológica de Chile INACAP \hfill 2013

\end{document}